\documentclass[11pt]{article}
\usepackage[margin=0.95in]{geometry}
\usepackage{amsmath,amssymb,amsthm}
\usepackage[T1]{fontenc}
\usepackage{lmodern}
\usepackage[hidelinks]{hyperref}
\usepackage{microtype}
\setlength{\parindent}{0pt}
\setlength{\parskip}{6pt}
\newtheorem{proposition}{Proposition}
\title{Failure of the increasing-inclusion claim in Conjecture 00000008586}
\author{gaochengzhecpu}
\date{October 4, 2026}
\begin{document}
\maketitle
\begin{abstract}
For the Hermitian matrix $T=\operatorname{diag}(0,1)$, the scalar $1/2$
belongs to the first compression range but not to the second. This
disproves the conjectured inclusion $\Lambda_1(T)\subseteq\Lambda_2(T)$.
The proof uses the rank-of-compression definition printed in the source.
The same matrix and scalar also contradict the proposed direction of
inclusion under the usual scalar-compression definition. The accompanying
Lean 4 project checks actual isometric compressions, their determinants,
and their matrix ranks.
\end{abstract}

\section{The source's definition}
Both language versions assert $\Lambda_k(T)\subseteq\Lambda_{k+1}(T)$.
They describe membership using the condition that the compression of
$\lambda I-T$ have rank at most $k-1$. Written with an orthonormal frame
for a $k$-dimensional subspace, this is
\begin{equation}\label{eq:source}
\begin{split}
\Lambda^{\rm src}_k(T)=\{\lambda\in\mathbb C:\ &\exists Q\in\mathbb C^{n\times k},\quad Q^*Q=I_k,\\[-2pt]
&\operatorname{rank}\bigl(Q^*(\lambda I_n-T)Q\bigr)\leq k-1\}.
\end{split}
\end{equation}
Here $Q^*$ is the conjugate transpose. The equation $Q^*Q=I_k$ says that
$Q$ is an isometric embedding of the standard complex Euclidean space
$\mathbb C^k$ into $\mathbb C^n$; the displayed product is its genuine
compression. We use \eqref{eq:source} for the main disproof.

\begin{proposition}
For $T=\left(\begin{smallmatrix}0&0\\0&1\end{smallmatrix}\right)$ and
$\lambda=1/2$, one has
$\lambda\in\Lambda^{\rm src}_1(T)$ but
$\lambda\notin\Lambda^{\rm src}_2(T)$.
\end{proposition}
\begin{proof}
Take the one-column frame
\[
 q=\frac1{\sqrt2}\begin{pmatrix}1\\1\end{pmatrix}.
 \qquad q^*q=1,\qquad q^*Tq=\frac12.
\]
Thus $q^*((1/2)I_2-T)q=0$ has rank zero, proving first-order membership.

For any two-column frame $Q$ with $Q^*Q=I_2$, taking determinants gives
$\det(Q^*)\det(Q)=1$. Put $B=(1/2)I_2-T=\operatorname{diag}(1/2,-1/2)$.
Then
\[
 \det(Q^*BQ)=\det(Q^*)\det(B)\det(Q)
            =\det(B)=-\frac14\ne0.
\]
Consequently every such compression has rank two, which exceeds
$2-1=1$. This excludes $1/2$ from $\Lambda^{\rm src}_2(T)$ and proves
the failure of the asserted inclusion.
\end{proof}

\newpage
\section{The usual scalar-compression convention}
The standard higher-rank numerical range is conventionally defined by
scalar compressions~\cite{CKZ}:
\[
 \Lambda^{\rm sc}_k(T)=\{\lambda\in\mathbb C:
   \exists Q\in\mathbb C^{n\times k},\ Q^*Q=I_k,\ Q^*TQ=\lambda I_k\}.
\]
This convention is kept separate from \eqref{eq:source}. For every such
scalar compression,
\[
 Q^*(\lambda I_n-T)Q
 =\lambda Q^*Q-Q^*TQ=0,
\]
so $\Lambda^{\rm sc}_k(T)\subseteq\Lambda^{\rm src}_k(T)$ for $k\geq1$.
The frame $q$ above gives $1/2\in\Lambda^{\rm sc}_1(T)$, while the
determinant argument excludes $1/2$ even from the larger source-defined
set at $k=2$. Thus the same witness disproves
$\Lambda^{\rm sc}_1(T)\subseteq\Lambda^{\rm sc}_2(T)$.

The inclusion is stated unconditionally. Our example uses the permitted
dimensions $n=2$, $k=1$, and $k+1=2$. It refutes this clause; the separate
generic-nonemptiness threshold, spectral envelopes, and arbitrary-dimensional
compactness are outside the delivered conclusion.

\section{Lean verification and reproducibility}
The Lean 4.19.0 project uses a pinned Mathlib revision and its actual
\texttt{Matrix.rank}, defined by the dimension of a linear-map range.
In \texttt{lean/Main.lean}:
\begin{itemize}
\item \texttt{statedRange} and \texttt{scalarRange} are the two distinct
existential sets above. The orthonormal-column equation is proved to induce
an actual Euclidean \texttt{Isometry} using the Hilbert-space adjoint.
\item The concrete frame identities are exact. For every two-column frame,
the determinant is $-1/4$; Mathlib's invertibility and rank theorems then
give rank two. The normalization uses the positive real square root.
\item The final theorems negate both inclusions and the universal
source-defined inclusion in dimension two, using the verified Hermitian
matrix and membership/exclusion pair.
\end{itemize}
Run these commands inside \texttt{lean/}:
\begin{verbatim}
lake build
lake env lean -DwarningAsError=true Main.lean
\end{verbatim}
The supplementary \texttt{verify.py} checks the witness using exact
fractions and expands the general determinant identity as a polynomial
in eight independent matrix-entry variables. Fresh-build logs, the axiom
audit, and hashes are in \texttt{verification/}. The bilingual source is
preserved byte for byte in \texttt{SOURCE.md}.

\small
\begin{thebibliography}{9}
\bibitem{Source} The Last Math Competition, conjecture 00000008586,
\href{https://github.com/The-Last-Math-Competition/The-Last-Math-Competition/blob/main/conjectures/00000008586.md}{repository statement}.
\bibitem{CKZ} M.-D. Choi, D. W. Kribs, and K. \.Zyczkowski,
\emph{Higher-Rank Numerical Ranges and Compression Problems},
Linear Algebra Appl. 418 (2006), 828--839.
\href{https://arxiv.org/abs/math/0511278}{arXiv:math/0511278}, definition (1).
\end{thebibliography}
\end{document}
